\bigskip\Needspace{10\baselineskip}
\pdfbookmark[1]{Fortgesetzte Nachrichten. Ceres Ferdinandea. März 1802}{ceres180203}
\begin{center}
{\small FORTGESETZTE NACHRICHTEN ÜBER DEN NEUEN HAUPTPLANETEN.\par}
CERES FERDINANDEA.\par
\medskip\hrule\medskip
Monatliche Correspondenz. Band V. S.\ 263—282. März 1802.\par
\medskip\hrule
\end{center}\medskip
—\quad —\quad —\quad „Ich kann nicht umhin, schreibt dieser verdienstvolle Gelehrte [\textsc{Gauss}], hier noch zu erwähnen, was für eine Wohltat es bei dieser Gelegenheit für die Astronomie war, dass es eine Zeitschrift wie die M.\ C.\ gab. Wie lau und gleichgültig hätte man die Entdeckung \textsc{Piazzi}s aufgenommen, wenn Sie nicht durch Ihre Zeitschrift alle Nachrichten darüber gesammelt, so schnell wie möglich verbreitet, das allgemeine Interesse geweckt, die Gründe dafür und dagegen abgewogen und es höchst wahrscheinlich gemacht hätten, dass dieses Gestirn ein Planet ist. Wahrscheinlich hätten sich nur wenige Astronomen die Mühe gemacht, es wieder zu suchen, da selbst der Lehrer und Meister aller heutigen Astronomen noch vor Kurzem so stark am neuen Planeten zweifelte?“ —\quad —\quad —

—\quad —\quad Im Dezemberheft des vorigen Jahres, S.\ 646 [S.\ 203 dieses Bandes], haben wir unseren Lesern alle Bahnelemente mitgeteilt, die Dr.\ \textsc{Gauss} nacheinander berechnet und anhand der Beobachtungen \textsc{Piazzi}s bis zum \textit{vierten Mal} verbessert hatte. An diesen \textit{vierten} Elementen hatte Dr.\ \textsc{Gauss} später noch die folgenden kleinen Änderungen vorgenommen:
\begin{quote}
$\AscendingNode\ 81^0\ 0'\ 13''39$, Aphel $326^0\ 27'\ 41''93$, Epoche $77^0\ 36'\ 38''35$, Neigung $10^0\ 37'\ 1''98$.
\end{quote}
Das Übrige blieb unverändert. —\quad —

—\quad —\quad —\quad Im vorangegangenen Februarheft, S.\ 178, haben wir bereits erwähnt, dass uns Dr.\ \textsc{Gauss} noch vor der Auffindung der \textit{Ceres} eine zum \textit{fünften Mal} verbesserte Bahn dieses Planeten zugeschickt hatte. Wir wollen daher auch diese Bahn vorerst zur historischen Kenntnis unserer astronomischen Leser gelangen lassen. Dr.\ \textsc{Gauss} berechnete die elliptische Bahn anhand der Beobachtungen \textsc{Piazzi}s wiederholt und mit großer Sorgfalt, wobei er Nutation, Aberration und Parallaxe aufs Genaueste berücksichtigte. Daraus ergaben sich die folgenden fünften Elemente:
\SourceLeaf{211}%
\begin{center}
(V)\par\smallskip
{\small\setlength{\tabcolsep}{2pt}
\begin{tabular}{l r}
Aphel $\ .\ .\ .\ .\ .\ .\ .\ .\ .$&$324^0\ 37'\ 11''$\\
$\AscendingNode\ .\ .\ .\ .\ .\ .\ .\ .\ .\ .\ .\ .$&$80\ 59\ 12$\\
Neigung $\ .\ .\ .\ .\ .\ .\ .\ .\ .\ .$&$10\ 37\ 9''55$\\
Logarithmus der großen Halbachse&$0.4446804$\\
Exzentrizität $\ .\ .\ .\ .\ .\ .\ .\ .\ .$&$0.0879111$\\
tägliche mittlere heliozentrische mittlere Bewegung&$763''950$\\
Epoche 1800 31. Dezember $\ .\ .\ .\ .$&$78^0\ 5'\ 16''6$
\end{tabular}}
\end{center}
Es wäre überflüssig, hier die schöne Übereinstimmung dieser Elemente mit den Beobachtungen \textsc{Piazzi}s zu zeigen, die Dr.\ \textsc{Gauss} aufs Sorgfältigste berechnet und bis auf ein paar Sekunden genau dargestellt hatte.

Als ich Dr.\ \textsc{Gauss} die Nachricht von der glücklichen Auffindung der so sehnlich erwarteten \textit{Ceres} und meine drei ersten Beobachtungen von ihr mitgeteilt hatte, berechnete er sie nach Erhalt als Erstes gleich nach seinen oben angegebenen V\textsuperscript{ten} Elementen. Er fand für die erste Beobachtung vom 7. Dezember in der Rektaszension einen Fehler von $+24'\ 8''$; für die zweite vom 11. Januar von $+30'\ 53''$; für die dritte vom 16. Januar von $+31'\ 53''$. Nach seinen IV\textsuperscript{ten} (verbesserten) Elementen weichen sie bei einer Überschlagsrechnung in derselben Reihenfolge so ab: $+14\frac12$ Minuten $+19\frac34$ Minuten $+20\frac12$ Minuten. Dass diese IV Elemente der Wahrheit etwas näher kommen als die V\textsuperscript{ten}, hält Dr.\ \textsc{Gauss} für Zufall; vielleicht liegt es aber auch zum Teil daran, wie sich die Störungen durch die Planeten auf die Beobachtungen \textsc{Piazzi}s auswirken, besonders auf die Breiten. Soweit Dr.\ \textsc{Gauss} aus den beiden ersten Beobachtungen von \textsc{Olbers} schließen konnte, die er inzwischen erhalten hatte, stimmten die Breiten mit beiden Sätzen von Elementen bis auf ein paar Minuten überein, und der Unterschied in den Deklinationen war hauptsächlich nur eine Folge der Unterschiede in der Länge. Bevor Dr.\ \textsc{Gauss} weitere Beobachtungen von mir erhielt, konnte er jedoch der Versuchung nicht widerstehen, seine Elemente vorläufig zu berichtigen. Dabei hielt er sich an meine beiden Rektaszensionen vom 7. Dezember und 16. Januar und fand so die folgenden VI\textsuperscript{ten} Elemente nach einer mit aller möglichen Sorgfalt ausgeführten Rechnung. Diese Sorgfalt war umso nötiger, als der geringste Fehler die ganze Arbeit nutzlos gemacht hätte.
\begin{center}
{\small\setlength{\tabcolsep}{2pt}
\begin{tabular}{l l}
Epoche 1801 Palermer Meridian $\ .\ .$&$77^0\ 24'\ 55''9$\\
—\quad1802\quad—\quad— $\ .\ .$&$155\ 33\ 35.1$\\
Tägliche mittlere tropische Bewegung&$770''7376$\\
Tropische Umlaufzeit $\ .\ .\ .\ .\ .\ .$&1681 Tage 12 St.\ 9 Min.\\
Logarithmus der großen Halbachse&$0.4421189$
\end{tabular}
\par\smallskip
$\left.\begin{array}{lr}
\text{Aphel}\ .\ .\ .\ .\ .\ .\ .\ .\ .&326^0\ 14'\ 45''\\
\AscendingNode\ .\ .\ .\ .\ .\ .\ .\ .\ .\ .\ .\ .&80\ 58\ 55
\end{array}\right\}\quad\begin{gathered}
\text{31. Dezember 1800,}\\
\text{siderisch ruhend angenommen}
\end{gathered}$
\par\smallskip
\begin{tabular}{l r}
Exzentrizität $\ .\ .\ .\ .\ .\ .\ .\ .\ .$&$0.08086253$\\
größte Mittelpunktsgleichung $\ .\ .\ .$&$9^0\ 16'\ 23''$\\
Neigung $\ .\ .\ .\ .\ .\ .\ .\ .\ .\ .\ .$&$10\ 37\ 51$
\end{tabular}}
\end{center}
Es lag in der Natur der bei dieser Rechnung verwendeten Methode, dass diese Elemente die Beobachtungen \textsc{Piazzi}s fast ebenso genau darstellen wie die V\textsuperscript{ten} Elemente. Es war jedoch nicht zu erwarten, dass gleich durch diese erste Verbesserung die Fehler bei den beiden neuen Beobachtungen sogleich von einem halben Grad auf null gebracht würden. Als Dr.\ \textsc{Gauss} inzwischen die beiden zugrunde gelegten Rektaszensionen nach diesen neuen Elementen berechnet hatte, fand er die Übereinstimmung doch größer, als er gehofft hatte. Der Fehler bei der ersten Beobachtung betrug nur $+3''0$, bei der letzten $+9''9$. Dr.\ \textsc{Gauss} hatte ferner drei neue Beobachtungen von mir, die ich ihm unterdessen geschickt hatte, nach diesen Elementen berechnet und fand die folgende Übereinstimmung in den Rektaszensionen.

\newcommand{\CeresSymbol}{\text{\mdseries\AstroSymbols ⚳}}
\SourceLeaf{212}%
\begin{center}
{\small\setlength{\tabcolsep}{5pt}
\begin{tabular}{r r|r|r|r}
\multicolumn{2}{c|}{Seeberg}&\shortstack{Berechnete\\Rektaszension $\CeresSymbol$}&Unterschied&\shortstack{Berechnete\\Deklination}\\\hline
1801 Dezember&7&$178^0\ 33'\ 33''6$&$+3''0$&$11^0\ 47'\ 33''$ N\\
1802 Januar&11&$186\ 46\ 9.3$&$+19.3$&$11\ 15\ 41$\\
—&; 16&$187\ 28\ 3.1$&$+9.9$&$11\ 26\ 40$\\
—&22&$188\ 6\ 45.9$&$+20.1$&$11\ 45\ 18$\\
—&25&$188\ 21\ 6.5$&$+27.3$&$11\ 56\ 49$
\end{tabular}}
\end{center}
Da bei der Bestimmung dieser Elemente nur die Breiten \textsc{Piazzi}s verwendet wurden und diese alle sehr nahe beim Knoten liegen, kann man für die Neigung noch keine große Genauigkeit erwarten. Nach einer Überschlagsrechnung mit den Deklinationen von \textsc{Olbers} glaubt Dr.\ \textsc{Gauss} jedoch, dass die Neigung kaum um mehr als eine Minute fehlerhaft sein dürfte. Das ist auch ein neuer Beleg für die ausgezeichnete Güte der Deklinationen \textsc{Piazzi}s.

Diese VI\textsuperscript{ten} Elemente werden später noch erhebliche Verbesserungen brauchen. Je weiter die Beobachtungen auseinanderliegen, desto sicherer fallen die Verbesserungen aus, verglichen mit einer möglichst genauen Anpassung der Elemente an die schon vorhandenen Beobachtungen durch Dr.\ \textsc{Gauss}. Dr.\ \textsc{Gauss} ist jedoch ein so geübter und unermüdlicher Rechner, dass ihm die Bestimmung neuer Elemente gar keine Mühe zu machen scheint. Ehe ich mich versah und ehe ich ihm meine Beobachtung vom 26. Januar zugeschickt hatte, überraschte er mich schon mit einer VII\textsuperscript{ten} Bahn der \textit{Ceres}. Dazu verwendete er anstelle meiner Beobachtung vom 16. die vom 25. Januar und erhielt damit die folgenden (VII) Elemente.
\begin{center}
(VII)
\par\smallskip
{\small
$\text{Epoche 1801}\ \left\{\begin{array}{lr}
\text{Seeburger Meridian}\ .\ .&77^0\ 27'\ 36''5\\
\text{Palermer Meridian}\ .\ .&77\ 27\ 30.9
\end{array}\right.$
\par\smallskip
\begin{tabular}{l r}
tägliche mittlere tropische Bewegung $\ .$&$769''7924$\\
Logarithmus der großen Halbachse $\ .$&$0.4424742$\\
Exzentrizität $\ .\ .\ .\ .\ .\ .\ .\ .\ .\ .$&$0.0814064$
\end{tabular}
\par\smallskip
$\left.\begin{array}{l}
\text{Aphel}\\
\AscendingNode
\end{array}\right\}\ \text{1801 ruhend}\ \begin{array}{lr}
.\ .\ .\ .\ .&325^0\ 57'\ 15''\\
.\ .\ .\ .\ .&80\ 58\ 40
\end{array}$
\par\smallskip
\begin{tabular}{l r}
größte Mittelpunktsgleichung $\ .\ .\ .$&$9\ 20\ 8.0$\\
Neigung $\ .\ .\ .\ .\ .\ .\ .\ .\ .\ .\ .$&$10\ 37\ 56.6$\\
{}[Handschrift] Epoche 1802 Palermo $\ .\ .$&$155\ 30\ 25.2$
\end{tabular}}
\end{center}
Hier wurden wiederum nur die Breiten \textsc{Piazzi}s verwendet. Diese Elemente stellen auch sämtliche Beobachtungen aus Palermo so genau wie möglich dar und stimmen jetzt auch besser mit meinen Beobachtungen überein, wie die folgende Übersicht zeigt:
\begin{center}
{\small\setlength{\tabcolsep}{4pt}
\begin{tabular}{r r|r|r|r|r}
\multicolumn{2}{c|}{Seeberg}&\shortstack{Berechnete\\Rektaszension}&Unterschied&\shortstack{Berechnete\\Deklination N.}&Unterschied\\\hline
1801 Dezember&7&$178^0\ 33'\ 29''2$&$-1''4$&&\\
1802 Januar&11&$186\ 45\ 47.6$&$-2.3$&&\\
&16&$187\ 27\ 38.8$&$-14.4$&&\\
&22&$188\ 6\ 18.2$&$-7.6$&&\\
&25&$188\ 20\ 37.2$&$-2.0$&$11^0\ 56'\ 58''4$&$+35''4$\\
&26&$188\ 24\ 37.0$&$-12.5$&$12\quad1\quad8.4$&\\
&28&$188\ 31\ 25.7$&$-12.1$&$12\quad9\ 55.6$&$+14.3$\\
&29&$188\ 34\ 14.1$&$-4.0$&$12\ 14\ 34.1$&\\
&30&$188\ 36\ 38.4$&$-5.5$&$12\ 19\ 19.8$&$+19.1$\\
&31&$188\ 38\ 38.3$&$-7.1$&$12\ 24\ 15.3$&\\
Februar&3&$188\ 42\quad9.5$&$-3.5$&$12\ 39\ 56.6$&$+17.6$\\
&4&$188\ 42\ 30.7$&$-5.6$&$12\ 45\ 25.7$&\\
&5&$188\ 42\ 26.7$&$-3.4$&$12\ 51\quad2.6$&$+34.7$\\
&9&$188\ 38\quad1.6$&$-2.3$&$13\ 14\ 43.0$&$+22.0$
\end{tabular}}
\end{center}
\expandafter\def\csname HistoricalNoteText18\endcsname{Damit wir auch hier nichts versäumen, was dazu beitragen kann, diesen Planeten möglichst lange und mit größtmöglicher Genauigkeit zu beobachten, werden wir jetzt, solange noch Zeit ist, alle Sterne vorab sehr genau zu bestimmen suchen, die in dem Teil der scheinbaren Bahn der \textit{Ceres} liegen werden, in dem man den Planeten nicht mehr im Meridian, sondern in seiner westlichen Quadratur nur mit Äquatorialinstrumenten oder parallaktischen Fernrohren beobachten kann. Diesen Beobachtungen ist auch deshalb größere Aufmerksamkeit zu widmen, weil sie für die Bestimmung des \textit{Radiusvektors} am günstigsten sein werden.}
\SourceLeaf{213}%
Es wäre für Dr.\ \textsc{Gauss} ein Leichtes, die oben gezeigte geringe Abweichung der Rechnung von allen bisher von mir beobachteten Rektaszensionen ganz verschwinden zu lassen. Da wir aber jetzt fast jede Woche von der Verbesserung dieser Elemente eine größere Genauigkeit erwarten können und sie auf jeden Fall hinreichend genau sind, um \textit{Ceres} noch lange aufzufinden, hält es Dr.\ \textsc{Gauss} selbst vorerst für das Beste, die Beobachtungen noch eine Zeit lang mit \textit{denselben Elementen} zu vergleichen und vor allem noch eine etwas größere Zahl genauer Deklinationen abzuwarten. Schon die geringe Abweichung, die die obige Tabelle zeigt, und die zu geringe Zahl zuverlässiger Deklinationen machen deutlich, dass es zur weiweitern Berichtigung der VII Elemente noch zu früh ist. Die Bahn der \textit{Ceres} ist ja wegen der Störungen durch die Einwirkung der übrigen Planeten keine vollkommene Ellipse. Diese Störungen müssen erst berechnet werden, ehe man die \textit{wahre} elliptische Bahn erhalten kann. —\quad —\quad —\quad Bei der jetzigen Lage, in der \textit{Ceres} bereits fast ein Viertel ihrer Bahn beschrieben hat, darf sich Dr.\ \textsc{Gauss} immerhin der Hoffnung hingeben, dass sich diese VII Elemente von den wahren Elementen der Bahn wohl nicht sehr stark unterscheiden. Mit allen Beobachtungen, auf die in diesem Jahr noch zu hoffen ist, bevor \textit{Ceres} unsichtbar wird\HistoricalNote{18}, hofft er diese Elemente so genau machen zu können, dass man sich damit bei einer Rückrechnung um 44 Jahre bis zu \textsc{Mayer}s Beobachtungen in der Länge um weniger als einen Grad und in der Breite nur um wenige Minuten irren kann. Sollte sich \textit{Ceres} anders als \textit{Uranus} in keinem älteren Sternverzeichnis finden, wird man ihre mittlere Bewegung freilich nur schrittweise genauer kennenlernen. Dann wird man noch lange jedes Jahr etwas hinzusetzen und verbessern müssen. —\quad —\quad —

Um unseren astronomischen Lesern auch künftig die Suche nach \textit{Ceres} zu erleichtern, teilen wir ihnen hier wieder eine kurze \textit{Ephemeride} ihres Laufs für die nächsten zwei Monate mit. Dr.\ \textsc{Gauss} hat sie nach seinen letzten und VII Elementen berechnet. Im April werden diese Stellungen hoffentlich um höchstens eine Minute vom Himmel abweichen.
\begin{center}
Stellungen der \textit{Ceres} im März und April 1802 für Mitternacht, mittlere Seeberger Zeit.
\par\smallskip
{\small\setlength{\tabcolsep}{5pt}
\begin{tabular}{r r|r|r|r}
\multicolumn{2}{c|}{1802}&\shortstack{Rektasz. $\CeresSymbol$\\in Grad}&Dekl.&Rektasz. in Zeit\\\hline
März&1&$186^0\ 41'$&$15^0\ 30'$ N&$12^{\mathrm u}\ 26^{\mathrm m}\ 45^{\mathrm s}$\\
&4&$186\ 11$&$15\ 50$&$12\ 24\ 45$\\
&7&$185\ 39$&$16\ 10$&$12\ 22\ 36$\\
&10&$185\quad5$&$16\ 29$&$12\ 20\ 18$\\
&13&$184\ 28$&$16\ 47$&$12\ 17\ 53$\\
&16&$183\ 51$&$17\quad4$&$12\ 15\ 24$\\
&19&$183\ 13$&$17\ 19$&$12\ 12\ 50$\\
&22&$182\ 34$&$17\ 33$&$12\ 10\ 15$\\
&25&$181\ 55$&$17\ 44$&$12\quad7\ 40$
\end{tabular}
\par\smallskip
\begin{tabular}{r r|r|r|r}
\multicolumn{2}{c|}{1802}&\shortstack{Rektasz. $\CeresSymbol$\\in Grad}&Dekl.&Rektasz. in Zeit\\\hline
März&28&$181^0\ 17'$&$17^0\ 54'$ N&$12^{\mathrm u}\quad5^{\mathrm m}\quad7^{\mathrm s}$\\
&31&$180\ 39$&$18\quad1$&$12\quad2\ 37$\\
April&3&$180\quad3$&$18\quad6$&$12\quad0\ 12$\\
&6&$179\ 29$&$18\quad9$&$11\ 57\ 54$\\
&9&$178\ 56$&$18\ 10$&$11\ 55\ 45$\\
&12&$178\ 26$&$18\quad9$&$11\ 53\ 46$\\
&15&$177\ 59$&$18\quad5$&$11\ 51\ 57$\\
&18&$177\ 35$&$17\ 59\frac12$&$11\ 50\ 20$
\end{tabular}}
\end{center}
Diesen Ephemeriden fügt Dr.\ \textsc{Gauss} noch die folgenden wichtigen Bemerkungen hinzu, die er daraus abgeleitet hat:

„Der Gegenschein des Planeten fällt auf den Nachmittag des 17. März. Um diese Zeit lässt er sich am günstigsten beobachten. Kurz vor der Opposition hat er den kleinsten Abstand von der Erde
\SourceLeaf{214}%
erreicht $=1.6025$. Ungefähr zur selben Zeit (etwas früher) erreicht er seine größte nördliche geozentrische Breite $=17^0\ 9'$, und etwas später seine schnellste rückläufige Bewegung mit täglich etwa 13 Min. in Rektaszension. Die nördliche Deklination nimmt noch bis Anfang April zu; ungefähr am 9. April wird der Lauf wieder südlich.“ —\quad —\quad —
\begin{flushright}[\textsc{von Zach}.]\end{flushright}

